% OWNER: sections/s2_4_feed.tex. Edit only this file.
\seclabel
\claim{\secnum{2.4}{sec:feed}\quad The live feed answers only when asked, and sends a delta only when the client echoes \code{t} exactly}

\lead{Connecting to the live feed takes 4 steps (\figref{signalr}).} Balances, orders and sector weights come through an ASP.NET Core SignalR hub at \code{/hubs/main}: an HTTP negotiation, a WebSocket upgrade and a handshake, after which the client sends feed requests and the server answers each one.

% OWNER: figs/signalr.tex. Edit only this file.
% Time-space diagram in the style of Kurose & Ross Fig 2.7/3.39: two timelines, time runs down, sloped arrows,
% message names along the arrows, events beside the timelines. The time axis is broken twice (zigzags).
\fig[wide]{signalr}{After negotiation and a handshake, the server answers only feed requests}[Messages between the client and the live feed's hub \code{/hubs/main}, time running down; set-up in milliseconds, then the 3\,s poll cycle]{%
\begin{tikzpicture}[x=1mm,y=1mm,
  f/.style={font=\fontsize{7}{8.4}\selectfont},
  host/.style={font=\fontsize{7.5}{9}\selectfont\bfseries,text=navy,anchor=south,align=center},
  tl/.style={greyM,line width=0.6pt},
  msg/.style={-{Stealth[length=1.6mm]},line width=0.7pt,navy},
  rep/.style={-{Stealth[length=1.6mm]},line width=0.7pt,greyD},
  ml/.style={font=\fontsize{7.5}{9}\selectfont,sloped,above,inner sep=0.7pt,text=navy},
  rl/.style={font=\fontsize{7.5}{9}\selectfont,sloped,below,inner sep=0.7pt,text=greyD},
  ev/.style={f,anchor=east,align=right,text=greyD},
  brk/.style={line width=0.6pt,greyM,fill=white},
  stp/.style={f,anchor=west,text width=48mm,align=left,text=greyD,execute at begin node={\hyphenpenalty=10000\exhyphenpenalty=10000\relax}},
  brace/.style={decorate,decoration={brace,amplitude=1mm},line width=0.5pt,navy}]
\def\C{0}\def\S{92}
% ---------------------------------------------------------------- hosts and timelines
\node[host] at (\C,1) {Client\\(browser or wrapper)};
\node[host] at (\S,1) {Server\\(SignalR hub)};
\draw[tl] (\C,0) -- (\C,-106); \draw[tl] (\S,0) -- (\S,-106);
\draw[tl,-{Stealth[length=1.4mm]}] (\C,-106) -- (\C,-109); \draw[tl,-{Stealth[length=1.4mm]}] (\S,-106) -- (\S,-109);
\node[f,text=greyD,anchor=north] at (\C,-109.3) {time};\node[f,text=greyD,anchor=north] at (\S,-109.3) {time};
% broken-axis marks: two zigzags across both timelines
\newcommand{\sigzig}[1]{\foreach \x in {\C,\S}{\fill[white] (\x-1.6,#1+0.9) rectangle (\x+1.6,#1-0.9);
  \draw[greyM,line width=0.5pt] (\x-1.6,#1+1.4) -- (\x+1.6,#1+0.4); \draw[greyM,line width=0.5pt] (\x-1.6,#1-0.4) -- (\x+1.6,#1-1.4);}}
% ---------------------------------------------------------------- set-up (milliseconds)
\draw[msg] (\C,-6) -- node[ml]{\code{POST /hubs/main/negotiate?negotiateVersion=1}} (\S,-10);
\draw[rep] (\S,-11) -- node[rl]{200: \code{connectionId}, \code{connectionToken}, 3 transports} (\C,-15);
\draw[msg] (\C,-23) -- node[ml]{\code{GET /hubs/main?id=...}, \code{Upgrade: websocket}} (\S,-27);
\draw[rep] (\S,-28) -- node[rl]{\code{101 Switching Protocols}} (\C,-32);
\draw[msg] (\C,-40) -- node[ml]{\code{{"protocol":"json","version":1}} 0x1E} (\S,-44);
\draw[rep] (\S,-45) -- node[rl]{\code{{}} 0x1E} (\C,-49);
% events beside the client timeline
\node[ev] at (\C-2,-8) {connection 1, HTTP/2,\\already open for REST};
\node[ev] at (\C-2,-25) {connection 2: new TLS\\connection, HTTP/1.1};
\node[ev] at (\C-2,-42) {handshake};
% ---------------------------------------------------------------- poll cycle
\draw[msg] (\C,-56) -- node[ml]{\code{RequestPortfolio} \code{{intv:3, auto:true, lastT:0, d, pfId}}} (\S,-60);
\draw[rep] (\S,-61) -- node[rl]{\code{ReceivePortfolio}: snapshot (\code{bs}, \code{ords}, \code{pos}, \code{msecs})} (\C,-65);
\draw[brace] (\C-2,-65) -- (\C-2,-56) node[midway,ev,xshift=-1.2mm] {reply in\\20 to 84\ms};
% 3 s wait with a break
\fill[navyM] (\C-0.6,-65.5) rectangle (\C+0.6,-76.5);
\sigzig{-71}
\draw[brace] (\C-2,-76.5) -- (\C-2,-65.5) node[midway,ev,xshift=-1.2mm] {wait 3\,s\\(not to scale)};
\draw[msg] (\C,-77) -- node[ml]{\code{RequestPortfolio} \code{{..., lastT = t of previous reply}}} (\S,-81);
\draw[rep] (\S,-82) -- node[rl]{\code{ReceivePortfolio}: delta (\code{pos_diff}, \code{ords_diff}, \code{msecs_diff})} (\C,-86);
\fill[navyM] (\C-0.6,-86.5) rectangle (\C+0.6,-93);
\sigzig{-90}
\node[ev] at (\C-2,-90) {wait 3\,s, then repeat};
% ping: both directions, drawn level (schematic)
\draw[{Stealth[length=1.6mm]}-{Stealth[length=1.6mm]},line width=0.6pt,greyD,dashed] (\C,-100) -- node[f,above,inner sep=0.8pt,text=greyD]{type 6 ping, sent by both sides; server every 15\,s} (\S,-100);
\node[ev] at (\C-2,-100) {keeps an idle\\connection open};
% ---------------------------------------------------------------- step brackets on the right
\foreach \a/\b/\n/\t in {-5/-16/1/{\textbf{Negotiate}: an HTTP/2 POST with the Bearer token on the REST connection; the reply offers WebSockets, Server-Sent Events and long polling},
  -22/-33/2/{\textbf{Connect}: the server offers no WebSockets over HTTP/2 [1], so the upgrade opens a second TLS connection. The wrapper sends the token as a header; browsers put it in the URL},
  -39/-50/3/{\textbf{Handshake}: both sides agree on JSON messages, each ended by the byte 0x1E},
  -55/-95/4/{\textbf{Poll}: the server sends a portfolio only in reply to a request. The client asks again 3\,s after each reply; a \code{lastT} equal to the previous reply's \code{t} returns a delta (\figref{lastt})}} {
  \draw[greyM,line width=0.5pt] (\S+3,\a) -- (\S+4,\a) -- (\S+4,\b) -- (\S+3,\b);
  \node[stp] at (\S+5.5,{(\a+\b)/2}) {\n.\ \t};
}
\end{tikzpicture}
}{Navy arrows are sent by the client and grey arrows by the server. Time runs down; the set-up part is drawn in milliseconds without measured durations, and the two zigzags mark the 3\,s waits, cut short. Reply times of 20 to 84\ms{} were measured before dawn with \code{auto} set to true and to false. The ping arrow is drawn level because both sides send pings independently. Type 1 is an invocation of a method on the other side, type 6 a ping and type 7 a close. In a browser recording from 28 Oct 2026 the client sent each request 3\,s after the previous reply.}{Browser developer tools, Python implementation, DART Quant Team}


\lead{The live feed sends nothing unless asked, before or during the session.} In the browser recording the browser asks again 3\,s after each reply. Our client got the same behaviour before dawn with the request flag \code{auto} set to true and to false (the browser sends true), with replies in 20 to 84\ms. From 09:05 to 09:08, with prices moving, a client that sent one feed request and then stayed silent for 60\,s received nothing but SignalR keep-alive pings (message type 6) every 15\,s.

\lead{On the same connection, echoing the previous reply's \code{t} returns only what changed.} A request with \code{lastT=0} returns a snapshot holding balances (\code{bs}), orders (\code{ords}), positions (\code{pos}) and market sector weights (\code{msecs}), plus 5 smaller top-level fields that Section~\ref{sec:guidelines} uses for the competition's guidelines. A request carrying the previous reply's \code{t} as \code{lastT} returns a delta in which those 4 fields are \code{null} and \code{pos_diff}, \code{ords_diff} and \code{msecs_diff} appear instead, each with the fields \code{added}, \code{removed}, \code{updated} and \code{ords_chgd}, together with \code{bs_updated}, a dict of the balance fields that changed. Right after our order (Section~\ref{sec:orders}) one reply had \code{null} arguments. We call this a null reply, as opposed to an empty delta, which has the delta fields but no changes; the receive loop has to accept a null reply and ask for a snapshot next.

% OWNER: figs/replysizes.tex. Edit only this file.
\fig{replysizes}{A delta is 1/17 of a snapshot before the open and about 1/8 during the session}[Size of one \code{ReceivePortfolio} reply in bytes, by reply type and time of day; 24 Nov 2026, polling every 3\,s]{%
{\tabfont\begin{tabularx}{\linewidth}{>{\RaggedRight\arraybackslash}X R{26mm} R{50mm} R{14mm}}
\thd{When} & \thd{Snapshot (\code{lastT=0})} & \thd{Delta (exact \code{t} echo)} & \thd{Ratio}\\
\rowsep
Before the open (before 09:00) & 9,915 & 577 (empty) & 17\\
\rowsep
During the session & 11,344 & 1,053 to 1,912 (mean about 1,400) & about 8\\
\end{tabularx}}
}{A snapshot holds balances (\code{bs}), orders (\code{ords}), positions (\code{pos}) and market sector weights (\code{msecs}) in full, plus 5 smaller top-level fields. In a delta those 4 fields are \code{null}; \code{pos_diff}, \code{ords_diff} and \code{msecs_diff} carry the records added, removed and updated, and \code{bs_updated} the changed balance fields. Before the open nothing changes, so the delta is empty. During the session \code{pos_diff.updated} listed 8 to 23 of the 23 holdings per reply, which sets the delta's range. Ratio = snapshot size divided by delta size (mean delta during the session).}{Python client on the live feed, DART Quant Team}


\lead{During the session a delta is about one eighth of a snapshot (\figref{replysizes}).} Before the open an empty delta was one seventeenth of a snapshot; with prices moving, deltas averaged about 1,400 bytes against a snapshot of 11,344. \code{pos_diff.updated} listed 8 to 23 of the 23 holdings per reply, and \code{msecs_diff} changed about every fourth reply (about 12\,s). \code{t} advanced by 3 per reply and sometimes by 4, so it follows wall-clock seconds.

{\RaggedRight \lead{\code{updated} carries only the changed fields of each record, keyed by the record's \code{Id} as a string.} Two entries from one reply read \code{"5745534": {"close": 89900, "prft": 1008000}} and \code{"5745535": {"close": 1180000, "prft": 3230000, "chPct": 8.76, "wei": 4.86}}: 2 to 4 of the 28 keys a position record has in a snapshot (Appendix~\ref{app:glossary} lists all 28). \code{ords_diff} has the same form, for example \code{{"added": [], "removed": [], "updated": {"1649882": {"cumQty": 6, "avgPx": 352416.67}}, "ords_chgd": {}}}. Snapshot records carry \code{Id} as an integer, so a delta key used as it arrives would match no record; the wrapper converts each key with \code{int()} before merging. The keys are position record IDs, and a ticker (\code{prodId}) is a separate field of the record, so the wrapper also keeps an index from \code{prodId} to \code{Id}.\par}

\lead{Merging deltas this way reproduced the server's positions and orders field for field in 15 of 15 rounds.} At 09:35, two connections polled the live feed at the same moments for 15 rounds. Connection A applied deltas, merging the \code{updated} fields by \code{int(Id)}; connection B asked for a full snapshot every round (\code{lastT=0}). In all 15 rounds both connections received the same \code{t}, and A's accumulated \code{pos} and \code{ords} records equalled B's snapshot field for field: 0 differences in 15 of 15 rounds. The check compared \code{pos} and \code{ords} only; balances and \code{msecs} were not compared field by field. Balances are kept current through \code{bs_updated}, which arrived as a dict in 20 of 20 deltas at 10:33 with the keys \code{prft}, \code{na}, \code{rt}, \code{nav} and \code{stkAmt}, sometimes also \code{pfStd} and \code{rank}; the wrapper merges it into \code{bs} (Section~\ref{sec:oneshare}). An earlier check compared the merged state with a snapshot taken about 3\,s after the last delta and found 13 records that differed only in \code{close}, \code{prft}, \code{wei} and \code{chPct}, the fields that move with prices in 3\,s; comparing replies with the same \code{t} removes that gap.

\lead{Of 3 candidate rules, only (a) fits all 11 replies: the server sends a delta only when \code{lastT} equals the \code{t} it last sent on this connection.} The rules were (a) \code{lastT} matches the \code{t} last sent on this connection; (b) \code{lastT} is any recent time; (c) \code{lastT} matches the \code{t} last sent to this account, on any connection. At 06:06, before the open, we varied \code{lastT} over 11 requests on 2 connections so that the rules predicted different replies (\figref{lastt}); every delta in that test was empty. The 09:35 merge check then ran during the session, when deltas carried changes, and connection A, echoing each \code{t}, kept receiving deltas whose merge matched the snapshot in all 15 rounds. That run is consistent with (a) but does not test (b) or (c), since all 3 rules predict a delta for an exact echo.

% OWNER: figs/lastt.tex. Edit only this file.
\fig[wide]{lastt}{Only rule (a) predicts every reply; (b) fails at 5 steps and (c) at 1}[Predicted and observed reply type by \code{lastT} sent; 11 requests 3\,s apart on 2 connections, 06:06 KST, 24 Nov 2026]{%
{\tabfont\begin{tabularx}{\linewidth}{R{6mm} >{\RaggedRight\arraybackslash}X >{\centering\arraybackslash}p{9mm} >{\centering\arraybackslash}p{17mm} >{\centering\arraybackslash}p{17mm} >{\centering\arraybackslash}p{17mm} >{\centering\arraybackslash}p{17mm}}
\thd{Step} & \thd{\code{lastT} sent} & \thd{Conn.} & \thd{(a) last \code{t}, this connection} & \thd{(b) any recent time} & \thd{(c) last \code{t}, this account} & \thd{Observed}\\
\rowsep
1 & 0 & A & \textcolor{greyD}{full} & \textcolor{greyD}{full} & \textcolor{greyD}{full} & \textbf{full}\\
\rowsep
2 & \code{t} of the previous reply & A & \textcolor{greyD}{delta} & \textcolor{greyD}{delta} & \textcolor{greyD}{delta} & \cellcolor{navyXL}\textbf{delta}\\
\rowsep
3 & \code{t} of the reply before that (6 s old) & A & \textcolor{greyD}{full} & \cellcolor{navyL}\textbf{delta\,$\times$} & \textcolor{greyD}{full} & \textbf{full}\\
\rowsep
4 & previous \code{t} minus 1 (never sent) & A & \textcolor{greyD}{full} & \cellcolor{navyL}\textbf{delta\,$\times$} & \textcolor{greyD}{full} & \textbf{full}\\
\rowsep
5 & previous \code{t} minus 60 & A & \textcolor{greyD}{full} & \cellcolor{navyL}\textbf{delta\,$\times$} & \textcolor{greyD}{full} & \textbf{full}\\
\rowsep
6 & previous \code{t} minus 3,600 & A & \textcolor{greyD}{full} & \textcolor{greyD}{n/a} & \textcolor{greyD}{full} & \textbf{full}\\
\rowsep
7 & previous \code{t} plus 100 & A & \textcolor{greyD}{full} & \textcolor{greyD}{n/a} & \textcolor{greyD}{full} & \textbf{full}\\
\rowsep
8 & \code{t} of the previous reply & A & \textcolor{greyD}{delta} & \textcolor{greyD}{delta} & \textcolor{greyD}{delta} & \cellcolor{navyXL}\textbf{delta}\\
\rowsep
8$'$ & the same value again & A & \textcolor{greyD}{full} & \cellcolor{navyL}\textbf{delta\,$\times$} & \textcolor{greyD}{full} & \textbf{full}\\
\specialrule{0.8pt}{1pt}{1pt}
9 & last \code{t} from connection A & B & \textcolor{greyD}{full} & \cellcolor{navyL}\textbf{delta\,$\times$} & \cellcolor{navyL}\textbf{delta\,$\times$} & \textbf{full}\\
\rowsep
10 & \code{t} of the previous reply & B & \textcolor{greyD}{delta} & \textcolor{greyD}{delta} & \textcolor{greyD}{delta} & \cellcolor{navyXL}\textbf{delta}\\
\end{tabularx}}
}{Mid-blue cells with $\times$ are predictions the observation contradicts; grey text marks predictions it confirms; in the Observed column, deltas are shaded light grey-blue; n/a: rule (b) does not say whether a \code{t} an hour old (step 6) or 100\,s ahead of the server (step 7) counts as recent, so it predicts neither reply type there. Step 3 sends a \code{t} the server really sent 6\,s earlier; step 4 a value it never sent; steps 5 to 7 test for a tolerance window; step 8$'$ repeats the value just used; step 9 opens connection B with connection A's last \code{t}. A reply is full when it contains \code{pos} and a delta when it contains \code{pos_diff}; before the open, full replies were 9,915 bytes and deltas 577 (\figref{replysizes}).}{DART Quant Team}


\lead{The server appears to keep, for each connection, the last snapshot it sent and its \code{t}.} When \code{lastT} equals that \code{t}, it sends the fields that changed since that snapshot; otherwise it sends everything. Either way, the current state becomes the new reference.

\lead{Rule (a) keeps the wrapper's state in step with the server, because a skipped delta always ends in a snapshot.} With one request in flight, a reply can be lost only as a timeout: the wrapper gives each request one deadline of 20\,s, which the server's pings every 15\,s do not extend, then drops the connection and reconnects, and a new connection starts with \code{lastT=0}, so the next reply is a snapshot and the client can never apply one delta after skipping another. A second request sent before the first reply arrives would carry a \code{lastT} that is already stale, so its reply would be a full snapshot, 17 times the size of an empty delta before the open (\figref{replysizes}).

\lead{\code{t} runs exactly 9 hours (32,400\,s) ahead of Unix time.} A reply received at 05:48:27 KST carried \code{t} = 1795499307. Read as UTC that is 05:48:27 on 24 Nov, matching the reply's \code{kstNow}, but the Unix time of 05:48:27 KST is 1795466907. The server evidently converts Korean wall-clock time to epoch seconds as if it were UTC. Other timestamps are correct: the JWT \code{exp} is true Unix time and order records carry UTC strings ending in \code{Z}. Under rule (a) the offset does no harm as long as the client echoes \code{t} without computing with it.

\lead{An open WebSocket keeps working after its token expires, but a new connection needs a valid token.} We opened a connection with a token expiring at 07:21:41 and kept polling every 3\,s (\figref{token}). Replies kept arriving for 8\,min past expiry, 382 in all, past the 300\,s after which REST requests are refused. A negotiation sent afterwards with the same token got a 401. The server did not recheck the token for 8\,min, which is consistent with ASP.NET Core SignalR's default of checking it only when a connection is made [3].

\lead{Any feed error makes the wrapper reconnect, and a feed silent for 26\,s blocks orders.} A 502, a 401, malformed JSON or a timeout closes the connection, and the wrapper reconnects after a random wait drawn uniformly from 0 to min(30, 2\textsuperscript{\textit{n}})\,s, where \textit{n} counts the failures in a row (full-jitter backoff). The wrapper records when the last reply arrived; reads, orders and \code{wait()} raise \code{StaleFeed} when that reply is older than 2\,×\,3 + 20 = 26\,s, 2 polling intervals plus the reply timeout. A 403 or 429 on the negotiation is the exception: the feed stops instead of reconnecting. Section~\ref{sec:design} describes how an operator sees and handles these signals.

\lead{The receive loop follows from these findings (\figref{loop}).} It keeps one request in flight, replaces its state on a snapshot, merges the changed fields of a delta into the records they name and \code{bs_updated} into \code{bs}, resets \code{lastT} to 0 after a null reply, and echoes \code{t} unread. On a reconnect it asks for a token, which renews it when due (Section~\ref{sec:auth}), negotiates again and starts from a snapshot.

% OWNER: figs/loop.tex. Edit only this file.
% Pseudocode of Hub.stream / poll_once / assert_fresh in pytimefolio/hub.py, comments beside the code.
\fig[wide]{loop}{The receive loop keeps one request in flight and reconnects after any error, starting again from a snapshot}[Pseudocode of the live-feed loop in \code{hub.py}, with the reason for each step beside it]{%
{\fontsize{7.6}{9.8}\selectfont
\newcommand{\cm}[1]{\leavevmode{\rmfamily\color{greyD}#1}}%
\newcommand{\I}[1]{\hspace*{#1mm}}%
\begin{tabular}{@{}>{\ttfamily\raggedright\arraybackslash}p{104mm} >{\RaggedRight\arraybackslash}p{70mm}@{}}
poll(hub):                                                   & \\
\I{4}n $\leftarrow$ 0; lastReply $\leftarrow$ never         & \cm{n counts consecutive failures}\\
\I{4}loop:                                                   & \\
\I{8}try:                                                    & \\
\I{12}if not connected:                                      & \\
\I{16}negotiate; open WebSocket; handshake                   & \cm{token renewed first if older than 720\,s or refused (Section~\ref{sec:auth})}\\
\I{16}lastT $\leftarrow$ 0                                   & \cm{a new connection starts from a snapshot}\\
\I{12}send(RequestPortfolio(lastT, d = currBizDay, pfId))   & \\
\I{12}p $\leftarrow$ await receive(timeout = 20\,s)          & \cm{one request in flight; pings answered while waiting}\\
\I{12}if p = null: lastT $\leftarrow$ 0                      & \cm{empty reply: state kept, next reply is a snapshot}\\
\I{12}elif p.pos $\neq$ null: state $\leftarrow$ p; lastT $\leftarrow$ p.t & \cm{snapshot: replace the state}\\
\I{12}else: merge(state, p); lastT $\leftarrow$ p.t          & \cm{delta: merge; \code{t} is echoed, never computed with}\\
\I{12}lastReply $\leftarrow$ now; n $\leftarrow$ 0           & \cm{recorded for the freshness check below}\\
\I{8}except any error:                                       & \cm{502, 401, malformed JSON, timeout, closed socket}\\
\I{12}close WebSocket; n $\leftarrow$ n + 1                  & \\
\I{12}await sleep(uniform(0, min(30, 2$^{\mathrm{n}}$)))       & \cm{full-jitter backoff, at most 30\,s; then reconnect}\\
\I{12}continue                                               & \\
\I{8}await sleep(3)                                          & \cm{ask again 3\,s after each reply}\\[3pt]
merge(state, p):                                             & \\
\I{4}for each diff in p.pos\_diff, p.ords\_diff, p.msecs\_diff: & \\
\I{8}add diff.added; drop diff.removed                       & \\
\I{8}for (i, f) in diff.updated: state[int(i)] $\mid$= f     & \cm{changed fields only; keys are record \code{Id}s as strings}\\
\I{4}state.bs $\mid$= p.bs\_updated                          & \cm{keeps net assets (\code{na}) current}\\[3pt]
assertFresh():                                               & \cm{called before reads, orders and \code{wait_done}}\\
\I{4}if now $-$ lastReply > 2 $\times$ 3 + 20 = 26\,s: raise StaleFeed & \cm{2 missed polls plus the 20\,s reply timeout}\\
\end{tabular}}}{\code{updated} keys are record \code{Id}s sent as strings; snapshot \code{Id}s are integers, hence \code{int(i)}. $\mid$= merges a dict of changed fields into the stored record. \code{uniform(0, x)} draws a wait uniformly between 0 and \code{x} seconds; a success resets n to 0. Token renewal is lazy: the feed asks for a token only when it connects or reconnects, and an open WebSocket keeps working after the token expires (\figref{token}). In a 15-round test against snapshots with the same \code{t} on a second connection, the merged \code{pos} and \code{ords} records matched field for field in every round (Section~\ref{sec:feed}).}{\code{pytimefolio/hub.py} v0.1.0, DART Quant Team}

\verdict{the receive loop keeps one request in flight and echoes \code{t} unread; when a time is needed, the wrapper uses the server clock taken from the replies' \code{Date} headers (Section~\ref{sec:auth}); it stores \code{kstNow} but does not use it (D8). Deltas are merged by \code{int(Id)}, changed fields only (D14). Token renewal leaves the open WebSocket alone; a reconnection asks for a token first, renewed when due, and sends it in a header, which keeps the access token out of URLs; the URL carries only the short-lived connection token (D9).}
